% ================================================================
%  ENCUENTRO 1 — Martes 28 de abril de 2026
% ================================================================
\chapter{Primer Encuentro}

\begin{destacado}
    \textbf{Fecha:} Martes, 28 de abril de 2026 \quad
    \textbf{Horario:} 6:15 a.m.\ -- 8:00 a.m. \quad
    \textbf{Asistentes:} Miguel, Kevin, Thomas, Isabella
\end{destacado}

\medskip

Este primer encuentro marcó el inicio formal del Semillero de Astronomía de la IEEOH.
Fue una sesión de presentación, diagnóstico y exploración colectiva de los intereses,
habilidades y expectativas de los integrantes. También se realizó el primer montaje del
telescopio como ejercicio introductorio.

% ---------------------------------------------------------------
\section{Diagnóstico Inicial}
% ---------------------------------------------------------------

Como primera actividad, se envió al grupo de WhatsApp del semillero un formulario de
diagnóstico individual con el propósito de conocer mejor a cada integrante y orientar
los próximos encuentros. El formulario incluye:

\begin{itemize}[leftmargin=1.5em]
    \item Grado escolar al que pertenece.
    \item Nivel de conocimiento previo en astronomía (escala del 1 al 5).
    \item Preguntas de selección múltiple sobre:
        \begin{itemize}
            \item ¿Qué aspecto de la astronomía te fascina más?
            \item ¿Cuál es tu objetivo personal dentro del semillero?
            \item ¿Qué habilidades traes que podrían aportar al grupo?
        \end{itemize}
\end{itemize}

\begin{actividad}{Formulario Diagnóstico}{act:form}
    El formulario está disponible en el siguiente enlace y debe ser diligenciado por
    todos los integrantes del semillero, incluidos quienes no asistieron al primer
    encuentro.

    \medskip
    \url{https://forms.gle/81MiR79YJZSrtt2A8}

    \medskip
    \textit{Nota: Los resultados serán socializados en el segundo encuentro.}
\end{actividad}

% ---------------------------------------------------------------
\section{Preguntas Orientadoras del Semillero}
% ---------------------------------------------------------------

En el primer encuentro se abrieron dos preguntas clave para explorar colectivamente
la identidad y el propósito del semillero. A continuación se recogen las respuestas
iniciales del grupo.

\subsection{¿Qué queremos hacer?}

Los integrantes manifestaron interés en las siguientes líneas de trabajo:

\begin{itemize}[leftmargin=1.5em]
    \item \textbf{Observación:} utilizar el telescopio para explorar el cielo nocturno
          y diurno; identificar planetas, la Luna, cúmulos estelares y otros objetos.
    \item \textbf{Entender teorías:} estudiar los fundamentos teóricos de la astronomía
          y la astrofísica (gravedad, luz, cosmología, sistema solar, etc.).
    \item \textbf{Hacer investigación:} desarrollar proyectos científicos propios,
          formular hipótesis y contrastarlas con datos reales.
    \item \textbf{Comunicar ciencia:} divulgar el trabajo del semillero a la comunidad
          escolar y al público en general por distintos medios.
    \item \textbf{Historia de la astronomía:} conocer el desarrollo histórico de la
          ciencia, sus protagonistas y los descubrimientos más relevantes.
\end{itemize}

\subsection{¿En qué puedo aportar?}

Se identificaron diversas formas en que los estudiantes pueden contribuir al semillero
según sus habilidades e intereses:

\begin{itemize}[leftmargin=1.5em]
    \item \textbf{Experimentación:} diseño y ejecución de actividades prácticas con el
          telescopio u otros instrumentos.
    \item \textbf{Divulgación:}
        \begin{itemize}
            \item Redes sociales (Instagram, TikTok).
            \item Página web del semillero.
            \item Emisora escolar.
            \item Periódico institucional.
        \end{itemize}
    \item \textbf{Fotografía:} registro visual de los encuentros y de las observaciones
          astronómicas.
    \item \textbf{Relaciones interdisciplinares:} generar vínculos con otros grupos del
          colegio que trabajen en arte, tecnología, música u otras áreas.
\end{itemize}

% ---------------------------------------------------------------
\section{Ideas y Proyectos Propuestos}
% ---------------------------------------------------------------

Durante la sesión emergieron diversas propuestas creativas. Se presentan aquí como
puntos de discusión y posibles proyectos futuros del semillero.

\begin{idea}{Sala oscura didáctica}{idea:salaoscura}
    Construir una sala oscura ambientada en astronomía que pueda funcionar como espacio
    de exposición tipo museo o planetario portátil. Los visitantes podrían recorrerla,
    aprender sobre el universo y tener una experiencia inmersiva del cielo nocturno.
    Este proyecto podría presentarse en ferias científicas o jornadas culturales.
\end{idea}

\begin{idea}{Teatro y narrativa astronómica}{idea:teatro}
    Crear una obra de teatro o performance artístico de temática astronómica: escenas
    sobre el Big Bang, el viaje de la luz, la vida de una estrella o personajes
    históricos como Copérnico, Galileo o Hypatia. Integra ciencia, historia y arte.
\end{idea}

\begin{idea}{Astronomía y música}{idea:musica}
    Explorar la relación entre la música y el cosmos: composiciones inspiradas en
    fenómenos astronómicos (como \textit{The Planets} de Holst), creación de piezas
    originales o uso de la sonificación de datos astronómicos reales.
\end{idea}

\begin{idea}{Salidas a espacios de astronomía}{idea:salidas}
    Planificar visitas a lugares especializados en la ciudad: el Planetario de Medellín,
    el Parque Explora, la Universidad de Antioquia (Instituto de Física) u observatorios
    privados de la Red de Astronomía de Colombia (RAC). También se consideran salidas
    nocturnas a zonas con baja contaminación lumínica para sesiones de observación.
\end{idea}

% ---------------------------------------------------------------
\section{Compromisos del Equipo}
% ---------------------------------------------------------------

En este primer encuentro cada integrante asumió un rol inicial dentro del semillero.
A continuación se describen los compromisos adquiridos y se proponen ideas concretas
y específicas para que cada uno pueda desarrollar su rol de manera efectiva.

\begin{compromiso}{Thomas — Registro fotográfico y redes sociales}{comp:thomas}
    \textbf{Rol asumido:} encargado del registro fotográfico de los encuentros y de la
    gestión de las redes sociales del semillero (Instagram y TikTok).

    \medskip
    \textbf{Ideas y propuestas concretas:}
    \begin{itemize}[leftmargin=1.5em]
        \item Crear y administrar las cuentas oficiales del semillero en Instagram y
              TikTok con un nombre y estética visual coherentes.
        \item Publicar semanalmente \textit{reels} o \textit{stories} sobre lo que se
              aprende en los encuentros: montaje del telescopio, curiosidades del cosmos,
              imágenes capturadas durante las observaciones.
        \item Lanzar retos o preguntas astronómicas interactivas para la audiencia
              (\textit{¿Sabes cuántas estrellas hay en la Vía Láctea?}, etc.).
        \item Documentar con fotografías el proceso completo de cada sesión: montaje
              del telescopio, reacciones de los estudiantes al observar por primera vez,
              experimentos realizados.
        \item Cubrir eventos astronómicos destacados (eclipses, lluvias de meteoritos,
              conjunciones) con publicaciones específicas preparadas con antelación.
        \item Proponer una estética visual coherente: paleta de colores, plantillas para
              publicaciones, hashtags propios del semillero.
    \end{itemize}
\end{compromiso}

\begin{compromiso}{Isabella — Periódico institucional}{comp:isabella}
    \textbf{Rol asumido:} responsable de la presencia del semillero en el periódico
    de la institución.

    \medskip
    \textbf{Ideas y propuestas concretas:}
    \begin{itemize}[leftmargin=1.5em]
        \item Escribir una columna fija mensual de astronomía en el periódico escolar:
              noticias recientes del cosmos, misiones espaciales en curso, efemérides
              del mes (lunas llenas, conjunciones, solsticios, etc.).
        \item Publicar perfiles de astrónomos y astrónomos históricos o contemporáneos,
              con énfasis en figuras latinoamericanas y mujeres en la ciencia.
        \item Redactar crónicas de los encuentros del semillero para mantener informada
              a la comunidad escolar sobre las actividades realizadas.
        \item Proponer una sección de preguntas y respuestas astronómicas donde los
              lectores envíen sus dudas y el semillero responda.
        \item Diseñar infografías o mapas estelares simples para incluir como
              suplemento imprimible en ediciones especiales del periódico.
        \item Elaborar entrevistas a los integrantes del semillero o a invitados
              externos (astrónomos, profesores universitarios, aficionados).
    \end{itemize}
\end{compromiso}

\begin{compromiso}{Miguel — Investigación}{comp:miguel}
    \textbf{Rol asumido:} líder de las tareas de investigación del semillero.

    \medskip
    \textbf{Ideas y propuestas concretas:}
    \begin{itemize}[leftmargin=1.5em]
        \item Proponer y liderar un primer mini-proyecto de investigación: por ejemplo,
              el seguimiento del movimiento de la Luna durante un mes (fases, posición
              en el cielo, relación con las mareas) con datos propios recopilados
              en las observaciones.
        \item Construir y mantener una \textit{bitácora científica del semillero}:
              registro sistemático de todas las observaciones (fecha, hora, objeto
              observado, condiciones atmosféricas, instrumento utilizado, descripción
              o boceto de lo observado).
        \item Elaborar un glosario astronómico progresivo: a medida que se estudian
              nuevos conceptos, documentar su definición, origen etimológico y contexto.
        \item Rastrear artículos científicos divulgativos (NASA, ESA, Ciencia NASA en
              Español) y preparar resúmenes breves para compartir con el grupo.
        \item Investigar las posibilidades de participar en la Olimpiada Colombiana
              de Astronomía u otras competencias científicas afines.
        \item Explorar bases de datos de acceso libre (como NASA ADS o Zooniverse)
              para participar en proyectos de \textit{citizen science}.
    \end{itemize}
\end{compromiso}

\begin{compromiso}{Kevin — Divulgación y apoyo transversal}{comp:kevin}
    \textbf{Rol asumido:} Kevin mostró disposición para aportar en múltiples áreas
    y mostró especial interés por la divulgación científica. Aún define su rol
    más específico.

    \medskip
    \textbf{Ideas y propuestas concretas:}
    \begin{itemize}[leftmargin=1.5em]
        \item Liderar la gestión de una \textit{página web} del semillero: espacio
              donde publicar artículos, bitácoras de observación, galerías de
              fotografías y recursos de aprendizaje para quienes quieran unirse.
        \item Preparar charlas o presentaciones cortas (\textit{talks} de 10--15 min)
              para divulgar temas astronómicos en otros cursos de la institución o en
              la emisora escolar.
        \item Explorar el rol de \textit{enlace} con otras áreas del colegio: proponer
              proyectos colaborativos con el grupo de tecnología, artes o ciencias
              sociales (por ejemplo, un mural artístico del sistema solar o un proyecto
              de física sobre movimiento orbital).
        \item Investigar y gestionar la relación con aliados externos: contactar a la
              Red de Astronomía de Colombia (RAC), a astrónomos aficionados locales o
              al Instituto de Física de la Universidad de Antioquia.
        \item Apoyar a Thomas en la producción de contenido para redes sociales cuando
              sea necesario, y colaborar con Isabella en la redacción de textos
              divulgativos.
    \end{itemize}
\end{compromiso}

% ---------------------------------------------------------------
\section{El Telescopio del Semillero}
% ---------------------------------------------------------------

En este primer encuentro se realizó el primer montaje del telescopio disponible en el
semillero. A continuación se presentan la descripción del instrumento y las reglas
fundamentales para su correcto uso y conservación.

\subsection{Descripción del Instrumento}

\begin{concepto}{Telescopio Decouvir Maksutov-Cassegrain MC 60/70}{con:telescopio}
    El telescopio del semillero es un \textbf{Decouvir Maksutov-Cassegrain MC~60/70},
    un instrumento de diseño catadióptrico (combina lentes y espejos) ampliamente
    valorado por su compacidad, portabilidad y alta nitidez de imagen.

    \medskip
    \textbf{Características principales:}
    \begin{itemize}[leftmargin=1.5em]
        \item \textbf{Diseño:} Maksutov-Cassegrain. Usa un espejo primario cóncavo,
              un espejo secundario convexo y una lente correctora menisco en la
              entrada del tubo.
        \item \textbf{Apertura (diámetro del objetivo):} 60\,mm.
        \item \textbf{Distancia focal:} 700\,mm (razón focal aproximada f/11.7).
        \item \textbf{Montaje:} sobre trípode con cabeza altacimutal.
        \item \textbf{Idóneo para observar:} Luna, planetas del sistema solar
              (Saturno y sus anillos, Júpiter y sus lunas galileanas), cúmulos
              abiertos y globulares, y objetos terrestres lejanos.
    \end{itemize}
\end{concepto}

\begin{nota}{Primer montaje del telescopio}{nota:montaje1}
    En este primer encuentro se realizó el montaje inicial del telescopio como
    ejercicio exploratorio. Por desconocimiento del procedimiento correcto, no
    fue posible enfocar ningún objeto. Queda pendiente consultar el manual y
    practicar el montaje sistemático y correcto en el próximo encuentro.
\end{nota}

\subsection{Reglas Fundamentales de Uso}

Las siguientes diez reglas son de obligatorio conocimiento y cumplimiento por parte de
\textbf{todos} los integrantes del semillero. Su propósito es proteger el instrumento,
garantizar observaciones exitosas y velar por la seguridad de los usuarios.
Las reglas técnicas más específicas y avanzadas se tratarán en una sección aparte a
medida que el grupo adquiera mayor experiencia.

\begin{regla}{No apuntar al Sol sin filtro solar certificado}{reg:sol}
    \textbf{Jamás} se debe apuntar el telescopio al Sol sin un filtro solar específico
    y certificado para uso astronómico. Observar el Sol directamente a través del
    telescopio, incluso por un instante, puede causar \textbf{ceguera permanente e
    irreversible}. Esta es la regla más importante y no admite excepciones.
    Antes de cualquier observación diurna, verificar que el filtro solar esté
    correctamente colocado y en buen estado.
\end{regla}

\begin{regla}{Aclimatar el telescopio antes de usarlo}{reg:aclimatar}
    Antes de comenzar una sesión de observación, sacar el telescopio del estuche o
    lugar de almacenamiento y dejarlo reposar al aire libre durante al menos
    \textbf{15 a 30 minutos}. Este proceso, llamado \textit{aclimatación térmica},
    permite que los elementos ópticos se equilibren con la temperatura ambiente y se
    eviten deformaciones en las imágenes causadas por el choque térmico.
\end{regla}

\begin{regla}{No tocar los elementos ópticos}{reg:optica}
    Los espejos y la lente correctora (menisco) del telescopio son extremadamente
    sensibles. \textbf{Nunca} se deben tocar con los dedos ni con ningún objeto
    que no sea el material de limpieza adecuado. La grasa natural de la piel
    deteriora los recubrimientos ópticos y genera manchas difíciles de eliminar.
    En caso de acumularse polvo, consultar el procedimiento de limpieza correcto
    antes de actuar.
\end{regla}

\begin{regla}{Montar el trípode de forma estable antes de colocar el tubo}{reg:tripode}
    Siempre se debe montar y nivelar el trípode \textbf{antes} de instalar el tubo
    óptico. Asegurarse de que las patas estén bien extendidas y fijas, y de que la
    cabeza esté nivelada. Un telescopio mal apoyado puede caer, lo que dañaría
    irreparablemente los espejos y la estructura del instrumento.
\end{regla}

\begin{regla}{Cubrir el objetivo y el ocular cuando no estén en uso}{reg:tapas}
    Las tapas protectoras del objetivo y del portaocular deben colocarse siempre que
    el telescopio no esté siendo utilizado activamente. Esto protege las ópticas
    del polvo, la humedad y la condensación. También se deben colocar las tapas de
    los oculares individuales cuando no estén montados en el telescopio.
\end{regla}

\begin{regla}{Guardar el telescopio en lugar seco, fresco y protegido}{reg:almacenamiento}
    El telescopio debe almacenarse en su estuche o funda protectora, en un lugar
    seco, fresco y sin exposición directa a la humedad, el polvo o variaciones
    térmicas extremas. La humedad es el principal enemigo de los instrumentos
    ópticos, ya que puede generar hongos sobre los espejos y la lente correctora.
    Se recomienda incluir bolsas de sílica gel dentro del estuche.
\end{regla}

\begin{regla}{No forzar ningún mecanismo de ajuste}{reg:mecanismos}
    Si algún ajuste, tornillo, rueda de enfoque o mecanismo del montaje ofrece
    resistencia, \textbf{no se debe forzar}. Antes de operar cualquier parte del
    telescopio, consultar el manual o preguntar a alguien con experiencia.
    Forzar mecanismos puede dañar el sistema de enfoque, la montura o el alineamiento
    interno del instrumento.
\end{regla}

\begin{regla}{Manipular el telescopio con cuidado durante el transporte}{reg:transporte}
    Al transportar el telescopio, hacerlo siempre dentro de su estuche protector
    y con el objetivo cubierto. No cargar el instrumento por el tubo óptico;
    utilizar el estuche o agarrar la montura. Evitar golpes, vibraciones fuertes
    y caídas. Un impacto puede desalinear los espejos y requerir una colimación
    profesional.
\end{regla}

\begin{regla}{No dejar el telescopio sin supervisión}{reg:supervision}
    El telescopio es un instrumento valioso y delicado. Cuando esté montado
    y en uso, debe haber siempre al menos una persona responsable a su cargo.
    Nunca dejarlo solo, especialmente en exteriores con viento, cerca de niños
    pequeños o en lugares con riesgo de volcarse.
\end{regla}

\begin{regla}{Registrar cada sesión de uso del telescopio}{reg:registro}
    Llevar un registro sistemático de cada vez que se usa el telescopio: fecha,
    lugar, condiciones atmosféricas (cielo despejado/nublado, humedad, viento),
    objetos observados y observaciones relevantes. Este hábito científico permite
    hacer seguimiento del estado del instrumento, detectar posibles problemas y
    construir una bitácora de observación que servirá como memoria del semillero.
\end{regla}

% ---------------------------------------------------------------
\section{Planes para el Próximo Encuentro}
% ---------------------------------------------------------------

Los siguientes puntos conforman la agenda tentativa del \textbf{segundo encuentro}:

\begin{enumerate}[leftmargin=1.5em, label=\arabic*.]
    \item \textbf{Socialización del primer encuentro} con los estudiantes que no
          pudieron asistir: presentar los compromisos asumidos, los acuerdos y los
          objetivos del semillero.
    \item \textbf{Presentación de resultados del formulario de diagnóstico:} analizar
          colectivamente las respuestas y extraer conclusiones sobre los intereses,
          niveles y expectativas del grupo.
    \item \textbf{Sesión enfocada en el telescopio:}
        \begin{itemize}
            \item Estudiar el \textit{Decouvir Maksutov-Cassegrain MC~60/70}:
                  nombre y función de cada parte, diseño óptico y especificaciones.
            \item Aprender el procedimiento correcto de montaje paso a paso.
            \item Revisar las reglas técnicas de uso.
            \item Realizar el primer montaje supervisado y correcto del telescopio.
            \item Primer ejercicio de observación: explorar objetos diurnos a distancia
                  (montañas, edificaciones lejanas) como práctica de enfoque antes de
                  pasar a objetos celestes.
        \end{itemize}
    \item \textbf{Asignación de tareas} para el siguiente periodo.
    \item \textbf{Planificación del tercer encuentro.}
\end{enumerate}

\begin{nota}{Pregunta abierta para el grupo}{nota:pregunta}
    ¿Qué podríamos observar en la mañana como primer ejercicio de observación
    astronómica/óptica? ¿Es posible usar el telescopio para observar montañas
    o construcciones lejanas como práctica de enfoque? Esta pregunta queda
    abierta para que los integrantes investiguen y lleguen con propuestas al
    siguiente encuentro.
\end{nota}
